\subsection{一个可分性判据}

命题 11.——设 K 的特征指数为 p，并令 E 为由集合 S 生成的 K 的代数扩张。若 E 在 K 上可分，则 \(\mathrm{E} = \mathrm{K}(S^{p^{n}})\) 对一切整数 \(n \geqslant 0\) 成立 ；反之，若 E 在 K 上次数有限且 \(\mathrm{E} = \mathrm{K}(S^{p})\) ，则 E 在 K 上可分。

情形 \(p = 1\) 由 V，第 37 页的推论可知是平凡的。以下设 \(p \neq 1\) 。

由假设，E 在 K 上是代数的，且我们有 \(E = K(S)\) ，由此可得

\[
\mathrm{K} \left(\mathrm{S} ^ {p}\right) = \mathrm{K} \left(\mathrm{E} ^ {p}\right) = \mathrm{K} \left[ \mathrm{E} ^ {p} \right] \quad \text { by   V,   p.18,   Cor.   1 }.
\]

若 E 在 K 上次数有限，则它是 K 的可分扩张当且仅当它是 K 上的平展代数；V，第 35 页的推论表明，这当且仅当 \(E = K[E^{p}]\) 时成立。

现设 E 在 K 上可分且次数无限。则 \(K[E^{p}]\) 是诸子环 \(K[E^{\prime p}]\) 的并，其中 \(E'\) 取遍 E 在 K 上次数有限的子扩张所成的集合；而这样的扩张 \(E'\) 在 K 上可分（V，第 36 页，命题 1），故由前所述有 \(E' = K[E^{\prime p}] \subset K[E^{p}]\) ；最后我们有 \(E = K[E^{p}]\) 。对 \(n \geqslant 0\) 作归纳，关系 \(E = K[E^{p}]\) 蕴含 \(E = K[E^{p^{n}}]\) 。

推论 1.——完全域的每个代数扩张都是完全域。

设 K 是特征指数为 p 的完全域，E 是 K 的一个代数扩张。则 E 在 K 上可分（V，第 36 页，命题 2），从而由命题 11 得 \(\mathrm{E} = \mathrm{K}(E^{p})\) ；但我们有 \(K = K^{p} \subset E^{p}\) ，因此 \(\mathrm{E} = \mathrm{K}(E^{p}) = E^{p}\) ，故 E 是完全域。

推论 2（Mac Lane）。——设 \(\bar{\mathbf{K}}\) 为 \(\mathbf{K}\) 的一个代数闭包，\(\mathrm{K}^{p^{-\infty}}\) 为 \(\mathbf{K}\) 在 \(\bar{\mathbf{K}}\) 中的完全闭包。要使子扩张 \(\mathbf{E}\) （\(\bar{\mathbf{K}}\) 的）在 \(\mathbf{K}\) 上可分，必要且充分的是它与 \(\mathrm{K}^{p^{-\infty}}\) 在 \(\mathbf{K}\) 上线性不相交。

我们可以立即归结为 \([E:K]\) 有限的情形。设 \((x_{i})_{i\in I}\) 为 E 在 K 上的一组基。E 与 \(K^{p^{-\infty}}\) 线性不相交，当且仅当它与 \(K^{p^{-n}}\) （对一切 \(n\geqslant0\) ）线性不相交，而这恰好意味着：关系式 \(\sum_{i\in I}x_{i}a_{i}^{p^{-n}}=0\) 蕴含 \(a_{i}=0\) （对一切 \(i\in I\) ），其中族

\((a_{i})_{i\in I}\) 由 K 的元素组成。这又意味着族 \((x_{i}^{p^{n}})_{i\in I}\) 在 K 上是自由的，或等价地说，它是 K 上向量空间 E 的一组基。换言之，E 与 \(\mathbf{K}^{p^{-n}}\) 线性不相交当且仅当 \(\mathrm{E} = \mathrm{K}(\mathrm{E}^{p^n})\) 。现在只需应用命题 11。

注记.—— 1) 当 E 在 K 上代数且次数无限时，条件 \(E = K(E^{p})\) 并不总能保证 E 在 K 上可分。例如，若 K 非完全而 E 是 K 的一个完全闭包，则我们有 \(E = K(E^{p})\) ，但 E 不是 K 的可分扩张（V，第 39 页，推论 3）。

\begin{enumerate}
\def\labelenumi{\arabic{enumi})}
\setcounter{enumi}{1}
\tightlist
\item
  设 \(\mathbf{E}\) 是域 \(\mathbf{K}\) （特征指数为 \(p\) ）的一个可分代数扩张。则我们有 \(\mathbf{E}^p \cap \mathbf{K} = \mathbf{K}^p\) （推论 2）；因此，若 \(\mathbf{E}\) 是完全域，则 \(\mathbf{K}\) 也是。
\end{enumerate}

